\section{Hints and solutions to the exercises}\label{sec:solutions}

\paragraph{Exercise~\ref{exo:reflgraph}.} Take objects $V$, $E$, arrows $s,t : V\to E$ and $r : E\to V$ with $r\circ s = r\circ t = \id_V$ (and the composites $s\circ r$, $t\circ r$ as further arrows $E\to E$). For a presheaf $F$, restriction along $r$ is a map $F(V)\to F(E)$ choosing a loop at each vertex, and $r\circ s = \id_V$ says that the source of that loop is the vertex itself; similarly for the target.

\paragraph{Exercise~\ref{exo:heyting}.} $A\cap\neg A = \emptyset$, so $A\subseteq\neg\neg A$ by maximality of $\neg\neg A$. Negation reverses inclusions ($A\subseteq B$ implies $\neg B\subseteq\neg A$); applying it to $A\subseteq\neg\neg A$ gives $\neg\neg\neg A\subseteq\neg A$, and $\neg A\subseteq\neg\neg(\neg A)$ is the first fact for $\neg A$. Finally $C$ is disjoint from $A\cup B$ iff it is disjoint from both, so $\neg(A\cup B) = \neg A\cap\neg B$.

\paragraph{Exercise~\ref{exo:dmlaw}.} If the De Morgan law holds for all $A,B$, take $B = \neg A$: $\neg(A\cap\neg A) = \neg\emptyset = F$, so $\neg A\cup\neg\neg A = F$. Conversely, assume $\neg A\cup\neg\neg A = F$. The inclusion $\neg A\cup\neg B\subseteq\neg(A\cap B)$ always holds. For the other, intersect $\neg(A\cap B)$ with $F = \neg A\cup\neg\neg A$: the part inside $\neg A$ is fine, and $\neg(A\cap B)\cap\neg\neg A$ is disjoint from $B$, because $\neg\neg A\cap B\subseteq\neg\neg A\cap\neg\neg B = \neg\neg(A\cap B)$, which is disjoint from $\neg(A\cap B)$. (The identity $\neg\neg A\cap\neg\neg B = \neg\neg(A\cap B)$ holds in every Heyting algebra.) So that part lies in $\neg B$.

\paragraph{Exercise~\ref{exo:three}.} $\mathbf 3$ has an initial object, so it satisfies Ore, and it is not a groupoid. Its Karoubi envelope is itself, so the essential points are the three evaluations. Its presheaf topos is $\Sh(X)$ for the Alexandrov space $X$ of the opposite chain, a three-point space whose open sets form a chain; its subtopoi correspond to its $2^3 = 8$ subsets (a finite $T_0$ space has one subtopos for each subset), and by Bridge~\ref{br:sub} to eight Grothendieck topologies on $\mathbf 3$.

\paragraph{Exercise~\ref{exo:oreproduct}.} Complete each component separately. In $\Delta\times\Delta$, the cospan $(\delta^1,\id_{[0]}), (\delta^0,\id_{[0]})$ into $([1],[0])$ cannot be completed, since its first components cannot. $\omega\times\omega$ is a product of chains.

\paragraph{Exercise~\ref{exo:karTn}.} If $a\in fMe$ and $b\in eMf$ with $ba = e$, $ab = f$, then on $\operatorname{im}e$ (where $e$ is the identity) $b\circ a$ is the identity, so $a$ is injective from $\operatorname{im}e$ into $\operatorname{im}f$; symmetrically, so the images have the same size. Conversely, for a bijection $\theta : \operatorname{im}e\to\operatorname{im}f$, put $a = \theta\circ e$ and $b = \theta^{-1}\circ f$. Every size from $1$ to $n$ occurs (fix $k$ points and send the others to one of them), and size $0$ does not.

\paragraph{Exercise~\ref{exo:automata}.} Write $R(s)$ for the set of states reachable from $s$; it is a trap. The weak excluded middle holds at $s$ for every trap iff any two states $t,t'\in R(s)$ have a common reachable state. (Necessity: take $A = R(t')$; then $s$ reaches $A$, so $t$ must reach $A$. Sufficiency: if $u\in R(s)\cap A$ and $t\in R(s)$, a common reachable state of $t$ and $u$ lies in $R(u)\subseteq A$.) So the automata for which it holds for all traps are those whose reachability preorder is confluent: this is Theorem~\ref{th:confluence} for that preorder. For instance, every automaton with a state reachable from all states qualifies, and the automaton of Figure~\ref{fig:automaton} does not.

\paragraph{Exercise~\ref{exo:finiteposets}.} By Theorem~\ref{th:confluence}, $[P,\Set]$ is De Morgan iff $P$ is confluent, and Boolean iff $P$ is discrete. Up to isomorphism, the posets with one or two elements are all confluent. With three elements there are five posets: the antichain, a two-element chain plus a point, the chain, the \q{V} ($a<b$, $a<c$) and the \q{$\Lambda$} ($b<a$, $c<a$). All except the V are confluent. Only the antichains give Boolean topoi.

\paragraph{Exercise~\ref{exo:semisimplicial}.} $\Delta_{\mathrm{inj}}$ fails the Ore condition, for the same reason as $\Delta$: the two vertices of $[1]$. So semi-simplicial sets are not De Morgan. With $\emptyset$ added, the category has an initial object, and the topos is De Morgan.

\paragraph{Exercise~\ref{exo:Nmult}.} The idempotents of $(\mathbb N,\cdot)$ are $0$ and $1$. The monoid is commutative, so they are not $\mathcal D$-related, and there are exactly two essential points: the forgetful functor $X\mapsto X$, and $X\mapsto X\cdot 0 = \{x : x\cdot 0 = x\}$, the \q{fixed points of $0$}. The topos is De Morgan (commutative, or: $0$ is a zero) and not Boolean ($2$ is not invertible).

\paragraph{Exercise~\ref{exo:inhabited}.} The sieve on $\emptyset$, in $\mathsf{Fin}^{\op}$, generated by the arrow that corresponds to the map $\emptyset\to\{\ast\}$ of $\mathsf{Fin}$: it says that the model with no generator must be covered by the model with one generator, i.e.\ that an element exists.
